% The Assertion Ladder: six rungs, each with its teaching move and its lean counterpart.
\begin{tikzpicture}[
    font=\small,
    rung/.style={draw=#1, line width=0.9pt, fill=#1!8, rounded corners=2pt,
                 minimum width=5.9cm, minimum height=0.78cm, anchor=west, align=left},
    note/.style={anchor=west, align=left, text width=4.6cm, font=\footnotesize},
  ]
  \def\gap{0.95}
  \foreach \i/\name/\teach/\lean/\col in {
      0/{L0 Implied}/{hidden curriculum}/{tacit knowledge, no standard}/muted,
      1/{L1 Shown}/{worked example}/{standard work, shown}/structure,
      2/{L2 Told}/{direct instruction}/{standard work, written}/requirement,
      3/{L3 Checked}/{assessment}/{built-in test at the point of work}/behaviour,
      4/{L4 Explained}/{formative feedback}/{andon: the stop says where and why}/reqc,
      5/{L5 Prevented}/{constraint}/{mistake-proofing (poka-yoke)}/structure} {
    \node[rung=\col] (r\i) at (0, \i*\gap) {\textbf{\name}\quad\textcolor{muted}{\footnotesize\teach}};
    \node[note] at (6.2, \i*\gap) {\textcolor{muted}{lean:} \lean};
  }
  % the rails of the ladder
  \draw[line width=1.6pt, jotblue] (-0.25,-0.5) -- (-0.25, 5*\gap+0.5);
  % arrows
  \draw[-{Stealth[length=3mm]}, line width=1pt, jotblue] (-0.8,-0.3) -- node[rotate=90, above, font=\footnotesize] {earlier, more certain feedback} (-0.8, 5*\gap+0.3);
  \draw[-{Stealth[length=3mm]}, line width=1pt, muted] (11.1, 5*\gap+0.3) -- node[rotate=-90, above, font=\footnotesize] {cheaper to build and maintain} (11.1,-0.3);
\end{tikzpicture}
